\documentclass[10pt,a4paper]{article}

\usepackage[utf8]{inputenc}
\usepackage[T1]{fontenc}
\usepackage{lmodern}
\input{glyphtounicode}
\pdfgentounicode=1
\usepackage[margin=0.52in]{geometry}
\usepackage[hidelinks]{hyperref}
\usepackage{enumitem}

\pagestyle{empty}
\setlength{\parindent}{0pt}
\setlength{\parskip}{0pt}
\setlist[itemize]{leftmargin=1.1em,itemsep=1pt,topsep=1pt,parsep=0pt}
\urlstyle{same}

\newcommand{\cvsection}[1]{%
    \vspace{0.4em}
    {\large\bfseries #1}\par
    \vspace{0.08em}
    \hrule
    \vspace{0.18em}
}

\newcommand{\jobentry}[4]{%
    \textbf{#1} \hfill \textbf{#2}\par
    \textit{#3}\par
    #4
    \vspace{0.15em}
}

\begin{document}

{\LARGE\bfseries ANDERSON H. R. FERREIRA}\par
\vspace{0.18em}
{\normalsize Data Scientist | Machine Learning Engineer | Doutorando Unicamp}\par
\vspace{0.28em}

Salto, SP | +55 11 95850-9690 | \href{mailto:a058899@dac.unicamp.br}{a058899@dac.unicamp.br}\par
\href{https://anderson-ferreira-83.github.io/}{Portfolio: anderson-ferreira-83.github.io} | \href{https://www.linkedin.com/in/anderson-ferreira-1a473138/?locale=pt}{LinkedIn: linkedin.com/in/anderson-ferreira-1a473138}\par
\href{https://github.com/anderson-ferreira-83}{GitHub: github.com/anderson-ferreira-83}\par
Defesa prevista ate dez/2026 | Disponibilidade hibrida em Sao Paulo entre jun/2026 e set/2026 | Elegivel ao Programa Academicos Itau 2026
\par Dados e Analytics | Linguagens de Programacao e Cloud | Machine Learning e IA | Matematica, Estatistica e Financas (base quantitativa transferivel) | Processos e Metodologia Agil

\cvsection{Resumo Profissional}
Doutorando na Unicamp em fase final de formacao, com atuacao em Dados e Analytics, Machine Learning e IA, Cloud e Estatistica. Experiencia em pipelines end-to-end, feature engineering, validacao de modelos e implementacao com FastAPI, Oracle e AWS. Combina pesquisa aplicada e execucao tecnica em projetos de IoT e saude, com tecnicas transferiveis para risco, fraude, credit scoring, churn e areas de negocio orientadas a dados.

\cvsection{Competencias-Chave}
Dados e Analytics: ETL, EDA, Feature Engineering, validacao de modelos. Linguagens e Cloud: Python, SQL, Java, JavaScript, FastAPI, Oracle SQL, AWS, Docker. Machine Learning e IA: scikit-learn, Random Forest, SMOTE, SHAP, threshold analysis. Matematica e Estatistica: Cross-validation, Holdout, analise estatistica. Agil: Scrum (fundamentos teoricos), backlog, sprint planning, review e retrospectiva.

\cvsection{Experiencia Profissional}
\jobentry{Professor Adjunto - Departamento de Computacao}{2023--Atual}{CEUNSP - Centro Universitario}{
\begin{itemize}
    \item Atuei como orientador e desenvolvedor em projeto de predicao de risco de hipertensao com 4.240 amostras e 12 variaveis clinicas.
    \item Implementei pipeline de treinamento com split estratificado, SMOTE apenas no treino e validacao cruzada 5-fold, preservando rigor metodologico e evitando data leakage.
    \item Conduzi avaliacao orientada a classes desbalanceadas, com foco em Recall, F2-score, AUC-ROC, thresholds clinicos e interpretabilidade com SHAP.
\end{itemize}
}

\jobentry{Pesquisador - Laboratorio de Vibroacustica}{2018--2026}{UNICAMP - Departamento de Mecanica Computacional}{
\begin{itemize}
    \item Publiquei pesquisa no MSSP 2026 (IF 8.9), comparando 5 geometrias com abordagem orientada a dados e validacao estatistica.
    \item Desenvolvi pipeline IoT end-to-end com ESP32 + MPU6050, FastAPI e Oracle XE, da coleta sensorial a classificacao de 7 estados operacionais.
    \item Expandi o espaco de atributos para 104 features candidatas e selecionei 16 finais; Random Forest de 200 arvores atingiu 97,34\% +- 0,63\% em validacao cruzada e 94,24\% em holdout.
\end{itemize}
}

\cvsection{Projetos Selecionados}
\textbf{IoT Anomaly Detection}\par
\begin{itemize}
    \item Projeto com mais de 210 mil registros em runs documentadas, 7 classes, 104 features candidatas e 16 finais para inferencia no browser sem servidor de ML dedicado.
    \item Mantive historico versionado de modelos e execucoes em CHANGELOG, MODEL\_REGISTRY e pipeline registry.
\end{itemize}

\textbf{Hyperten ML}\par
\begin{itemize}
    \item Projeto de predicao de risco de hipertensao com 4.240 amostras e 12 variaveis, usando SMOTE apenas no treino para evitar data leakage em problema desbalanceado.
    \item Avaliacao orientada a Recall e F2-score, com thresholds clinicos, SHAP e arquitetura AWS baseada em S3, CloudFront, API Gateway e Lambda.
\end{itemize}

\cvsection{Formacao Academica}
\textbf{Doutorado em Engenharia Mecanica} - UNICAMP \hfill 2018--2026\par
Defesa prevista ate dez/2026\par
\textbf{Mestrado em Engenharia Mecanica} - UNICAMP \hfill 2011--2013\par
\textbf{Licenciatura em Fisica} - UNICAMP \hfill 2006--2010

\cvsection{Certificacoes e Publicacao}
AWS AI Practitioner (2025) | Applied ML with Python (2024) | IBM Data Engineer Professional (2023) | GCP Professional Cloud Architect (2023)\par
Mechanical Systems and Signal Processing (MSSP), 2026 | CiteScore 18.0 | Impact Factor 8.9

\end{document}
